\documentclass[10pt,a4paper]{amsart}
\usepackage[margin=19mm]{geometry}
\usepackage{amsmath,amssymb,mathtools}
\usepackage[scr=rsfs,cal=euler]{mathalfa}
\usepackage{xfrac}
\usepackage[unicode,hidelinks,bookmarks=false]{hyperref}
\newcommand{\ie}{i.e.\ }
\newcommand{\cf}{cf.\ }
\newcommand{\resp}{respectively\ }
\newcommand{\Subsection}{\subsection}
\providecommand{\nobreakdash}{\nobreak\hbox{-}}
\setlength{\emergencystretch}{2em}
\multlinegap0pt
\usepackage{amssymb}
\usepackage{mathtools}
\usepackage{enumerate}
\usepackage{float}
\newtheorem{prop}{Proposition}
\title{Symmetric spaces for groups over involutive algebras and applications to Higgs bundles}
\author{Pengfei Huang, Georgios Kydonakis, Eugen Rogozinnikov, Anna Wienhard}
\date{Source: arXiv:2509.02002v2 (CC BY 4.0)}
\begin{document}
\maketitle

\noindent\emph{Source and licence.} Symmetric spaces for groups over involutive algebras and applications to Higgs bundles. Version arXiv:2509.02002v2 (CC BY 4.0), distributed by arXiv under the Creative Commons Attribution 4.0 International licence, http://creativecommons.org/licenses/by/4.0/, as recorded in the arXiv licence metadata for this exact version and shown on the abstract page https://arxiv.org/abs/2509.02002v2. Full licence notice: \texttt{licenses/arxiv-2509.02002-CC-BY-4.0.txt}.\par\smallskip

\noindent\textit{Editorial scope.} Counted unit: the article's prop prop:incarn (source0.tex lines 852--870). The article's complete original proof of it is delivered with it (source0.tex lines 872--923, one \texttt{proof} environment of 52 lines). No mathematics is written, repaired or re-derived: the statement and the proof are the article's own lines, in the article's own notation, and no retained line is substituted or re-flowed.  No further source-local context is needed: every cross-reference of every retained line resolves inside the two delivered spans.  Nothing was omitted from any retained span, so the package carries no omission record.  No retained line contains a citation, so the wrapper carries no bibliography and no bibliography entry.  No retained line contains a figure environment, an image-inclusion command or drawing code, so no image is delivered and the package declares no asset dependency.  Editorial formatting only, in addition: an AMS \texttt{amsart} preamble that restates the article's own declarations and macros used by the retained lines, namely the environments \texttt{prop} on separate counters, together with the standard packages those lines need.  The retained environments print in this document as Proposition 1 in order of appearance, whereas the article prints the same items with the numbering of its own full document.  The complete native source of the article as distributed in the arXiv e-print for this exact version is bundled unchanged as \texttt{sources/184290/source0.tex}.
\begin{prop}\label{prop:incarn}
    Let $(A,\sigma)$ be a real involutive algebra. Then, the following assertions hold:
    \begin{enumerate}    
        \item The group $\mathrm{O}_{(1,1)}(A_\mathbb C,\bar\sigma_\mathbb C)$ is isomorphic to 
        $\mathrm{Sp}_{2}(A_\mathbb C,\bar\sigma_\mathbb C)$. Moreover,
        \begin{align*}
        \mathrm{Sp}_2(A,\sigma)&=\mathrm{Sp}_{2}(A_\mathbb C,\sigma_\mathbb C)\cap \mathrm{Sp}_{2}(A_\mathbb C,\bar\sigma_\mathbb C), \text{ and }\\
         T^{-1}\mathrm{Sp}_2(A,\sigma)T&=\mathrm{Sp}_{2}(A_\mathbb C,\sigma_\mathbb C)\cap R^{-1}\mathrm{O}_{(1,1)}(A_\mathbb C,\bar\sigma_\mathbb C)R,
        \end{align*}
        where $T=\frac{1}{\sqrt{2}}\begin{pmatrix}
        1 & i \\
        i & 1 \end{pmatrix}$ and $R=\frac{1}{\sqrt{2}}\begin{pmatrix}
        1 & 1 \\
        -1 & 1 \end{pmatrix}$.
        \item The group $\mathrm{O}_{(1,1)}(A_\mathbb H,\sigma_1)$ is isomorphic to $\mathrm{Sp}_{2}(A_\mathbb H,\sigma_0)$.
        \item The group $\mathrm{O}_{(1,1)}(A_\mathbb H,\sigma_0)$ is isomorphic to $\mathrm{Sp}_{2}(A_\mathbb H,\sigma_1)$.
        \item $\mathrm{Sp}_{2}(A_{\mathbb C\{i\}},\sigma_\mathbb C)=\mathrm{Sp}_{2}(A_\mathbb H,\sigma_0)\cap \mathrm{Sp}_{2}(A_\mathbb H,\sigma_3)$, where $\sigma_3(x+yj)={\sigma}_{\mathbb{C}}(x)-\bar{\sigma}_{\mathbb{C}}(y)j$, for $x,y\in A_{\mathbb C\{i\}}$.    
    \end{enumerate}
\end{prop}

\begin{proof}
(1). In the basis $e_1:=(1,0)^t$, $e_2:=(0,i)^t$ of $A_\mathbb C^2$, the standard indefinite form $\omega$ is represented by the matrix $\Omega=\begin{pmatrix}
    0 & i \\
    -i & 0
\end{pmatrix}$. This change of basis conjugates the group $\mathrm{O}_{(1,1)}(A_\mathbb C,\bar\sigma_\mathbb C)$ by the matrix $S=\frac{1}{\sqrt{2}}\begin{pmatrix}
    1 & 0 \\
    0 & i
\end{pmatrix}$.
The form $\Omega$ is proportional to the standard symplectic form. Therefore, the conjugated group $S^{-1}\mathrm{O}_{(1,1)}(A_\mathbb C,\bar\sigma_\mathbb C)S$ agrees with $\mathrm{Sp}_{2}(A_\mathbb C,\bar\sigma_\mathbb C)$.  Moreover, if $g\in \mathrm{Sp}_{2}(A_\mathbb C,\bar\sigma_\mathbb C)\cap \mathrm{Sp}_{2}(A_\mathbb C,\sigma_\mathbb C)$, then $\sigma_\mathbb C(g)=\bar\sigma_\mathbb C(g)$, i.e. $g\in\mathrm{Mat}_2(A)$.
\vspace{2mm}

\noindent (2)--(4). We consider the following indefinite form on $A_{\mathbb{H}}^{2}$:
\[h(x,y):={{\sigma }_{1}}{{(x)}^{t}}\begin{pmatrix}
   0 & j  \\
   -j & 0  \\
\end{pmatrix}y.\]
We now check that the form $h$ is ${{\sigma }_{1}}$--sesquilinear. Indeed, one sees: 
\begin{align*}
h(y,x)  & =\sigma_1(y)^t\begin{pmatrix}
   0 & j  \\
   -j & 0  \\
\end{pmatrix}x\\
 & =\sigma_1\left( \sigma_1(x)^t\begin{pmatrix}
   0 & j  \\
   -j & 0  \\
\end{pmatrix}y \right)\\
& =\sigma_1\left( h(x,y) \right).
\end{align*}
Then, with respect to the basis $e_1:=\left(1,0\right)^t$, $e_2:=\left(0,-j\right)^t$, we have $h\left( {{e}_{1}},{{e}_{1}} \right)=h\left( {{e}_{2}},{{e}_{2}} \right)=0$, and $h\left( {{e}_{1}},{{e}_{2}} \right)=1$, that is, $h$ is indefinite. Thus, 
$$S^{-1}\mathrm{Aut}(h)S=\mathrm{O}_{(1,1)}(A_\mathbb H,\sigma_1),$$
for $S=\begin{pmatrix}
    1 & 0 \\
    0 & -j 
\end{pmatrix}$.

 \noindent Further, $h=j\hat\omega_\mathbb H$, where $\hat\omega_\mathbb H(x,y)=\sigma_3(x)\begin{pmatrix}
   0 & 1  \\
   -1 & 0  \\
\end{pmatrix} y$ where $\sigma_3(x+yj)={\sigma}_{\mathbb{C}}(x)-\bar{\sigma}_{\mathbb{C}}(y)j$. This means that $\mathrm{Sp}_2(A_\mathbb H,\sigma_3)=\mathrm{Aut}(h)$. Finally, the involution $\theta_i\colon \mathbb H\to\mathbb H$, $q\mapsto i^{-1}qi$, for $q\in\mathbb H$ fixing $\mathbb C\{i\}\subset\mathbb H$, induces an~involution $\theta_j\colon\mathrm{Mat}_2(A_\mathbb H)\to \mathrm{Mat}_2(A_\mathbb H)$ fixing $\mathrm{Mat}_2(A_{\mathbb C\{i\}})$ such that $\theta_j\circ\sigma_3\circ\theta_j=\sigma_0$. Thus, 
$$
\theta_j\mathrm{Sp}_2(A_\mathbb H,\sigma_3)\theta_j=\mathrm{Sp}_2(A_\mathbb H,\sigma_0).$$

\noindent  Similarly, because $(ij)^{-1}\sigma_0(g)ij=\sigma_1(g)$, we obtain $\mathrm{Sp}_2(A_\mathbb H,\sigma_1)=\mathrm{Aut}(\hat h)$, where 
$$\hat h(x,y)=\sigma_0(x)^t\begin{pmatrix}
    0 & ij \\
    -ij & 0
\end{pmatrix}y.$$
In the basis $e_1=(1,0)^t,\;e_2=(0,ij)^t$, this form agrees with the standard indefinite form.

 \noindent Finally, if $g\in \mathrm{Sp}_{2}(A_\mathbb H,\sigma_0)\cap \mathrm{Sp}_{2}(A_\mathbb H,\sigma_3)$, then $\sigma_3(g)=\sigma_0(g)$, i.e. $g\in\mathrm{Mat}_2(A_{\mathbb C\{i\}})$.

\end{proof}

\end{document}
